\documentclass[11pt,a4paper]{article}

\usepackage{amsmath,amssymb}
\usepackage{booktabs}
\usepackage{longtable}
\usepackage[margin=1.0in]{geometry}
\usepackage{hyperref}
\usepackage{fancyhdr}

\pagestyle{fancy}
\fancyhf{}
\fancyhead[R]{\small\itshape FTD Constants Reference}
\fancyfoot[C]{\thepage}
\renewcommand{\headrulewidth}{0.4pt}

\newcommand{\Gstar}{{G^{*}}}
\newcommand{\vpi}{\varpi}

\title{\textsc{FTD Constants Reference}\\[6pt]
\large Canonical Values to 25-Digit Precision\\
Computed from $D = 3$ and $\vpi$}

\author{}
\date{}

\begin{document}
\maketitle
\thispagestyle{empty}

\begin{quote}\itshape
All dimensionless quantities are computed from $D = 3$ and $\vpi$; one external unit calibration $K_B = 0.5100$~MeV sets SI units. The Planck mass $m_P = K_B / (\sqrt{2\pi}\cdot\tfrac{16}{3}\cdot\alpha^{11})$ is derived, not an independent input.
\end{quote}

\smallskip

\noindent\textbf{Conventions.} $D = 3$ is the spatial dimension (axiom); $\mathcal{D} = N_c N_{\mathrm{base}}^2 - 1 = 47$ is the constraint denominator (derived integer). All masses and couplings use the 4-term corrected $\alpha$ unless tagged ``tree.'' The tree-level value $\alpha_{\mathrm{tree}} = 1/x_+$ appears only in the Layer~5 table and in $r_{\mathrm{eff}}$, where the uncorrected root is the relevant quantity. $T(n) = n(n+1)/2$ is the $n$th triangular number.

\bigskip

%=======================================================================
\section*{Layer 0: Transcendental Seed}
%=======================================================================

\begin{longtable}{lp{5.5cm}r}
\toprule
\textbf{Symbol} & \textbf{Formula} & \textbf{Value} \\
\midrule
\endhead
$\Gamma(1/4)$ & Euler gamma function at $1/4$ &
$3.6256099082219083119306851\ldots$ \\
\bottomrule
\end{longtable}

%=======================================================================
\section*{Layer 1: Lemniscatic Geometry}
%=======================================================================

\begin{longtable}{lp{5.5cm}r}
\toprule
\textbf{Symbol} & \textbf{Formula} & \textbf{Value} \\
\midrule
\endhead
$\vpi$ & $\displaystyle 2\!\int_0^1 \frac{dx}{\sqrt{1-x^4}}$%
\footnote{Closed form: $\vpi = \Gamma(1/4)^2/(2\sqrt{2\pi})$. The integral definition is $\pi$-free.} &
$2.6220575542921198104648396\ldots$ \\[8pt]
$M$ & $\displaystyle\frac{\vpi}{\pi} = \frac{1}{\mathrm{AGM}(1,\sqrt{2})}$ &
$0.8346268416740731862814297\ldots$ \\
\bottomrule
\end{longtable}

%=======================================================================
\section*{Layer 2: Bridge Constant}
%=======================================================================

\begin{longtable}{lp{5.5cm}r}
\toprule
\textbf{Symbol} & \textbf{Formula} & \textbf{Value} \\
\midrule
\endhead
$\Gstar$ & $\displaystyle\frac{\Gamma(1/4)^2}{\sqrt{2}\;\Gamma(1/2)^2}$%
\footnote{Equivalent forms: $\Gstar = 2\vpi/\sqrt{\pi} = \sqrt{2}\,\Gamma(1/4)^2/(2\pi)$. The $\Gamma$-ratio form is $\pi$-free: both $\Gamma(1/4)$ and $\Gamma(1/2)$ are definable by integrals involving only $e^{-t}$ and rational powers of~$t$. The identity $\Gamma(1/2)^2 = \pi$ is a consequence, not an input.} &
$2.9586751191886388923108214\ldots$ \\[8pt]
$\pi$ & $\displaystyle\frac{4\,\vpi^2}{\Gstar^2} = 4\ell$ \quad (derived: not circular\footnote{$\vpi$ is computed by the $\pi$-free integral $2\!\int_0^1 dx/\sqrt{1-x^4}$. $\Gstar$ is computed by the $\pi$-free $\Gamma$-ratio above. Their combination gives $\pi = 4\vpi^2/\Gstar^2$ without $\pi$ appearing in any input.}) &
$3.1415926535897932384626434\ldots$ \\[8pt]
$\ell$ & $\displaystyle\frac{\vpi^2}{\Gstar^2} = \frac{\pi}{4}$ &
$0.7853981633974483096156608\ldots$ \\
\bottomrule
\end{longtable}

%=======================================================================
\section*{Layer 3: Master Quadratic}
%=======================================================================

\noindent Equation: $\;x^2 - 16\,\Gstar^2\, x + 16\,\Gstar^3 = 0$

\begin{longtable}{lp{5.5cm}r}
\toprule
\textbf{Symbol} & \textbf{Formula} & \textbf{Value} \\
\midrule
\endhead
$x_+$ & $8\Gstar^2\!\left(1 + \sqrt{1 - \frac{1}{4\Gstar}}\right)$ &
$137.03617145815548388160057\ldots$ \\[8pt]
$x_-$ & $8\Gstar^2\!\left(1 - \sqrt{1 - \frac{1}{4\Gstar}}\right)$ &
$3.0239639163390210039527059\ldots$ \\[8pt]
$x_+\! \cdot\! x_-$ & $16\,\Gstar^3$ &
$414.39243772270942760105875\ldots$ \\[4pt]
$x_+\!+\!x_-$ & $16\,\Gstar^2$ &
$140.06013537449450488555328\ldots$ \\[4pt]
$\Delta$ & $64\,\Gstar^3(4\Gstar - 1)$ &
$17959.271770230889251899402\ldots$ \\
\bottomrule
\end{longtable}

%=======================================================================
\section*{Layer 3b: Precision Formula}
%=======================================================================

\noindent $\displaystyle\alpha^{-1} = x_+ - c_1|\varepsilon| + c_2|\varepsilon|^2 - c_3|\varepsilon|^3 - c_4|\varepsilon|^4$

\begin{longtable}{lp{5.5cm}r}
\toprule
\textbf{Symbol} & \textbf{Formula} & \textbf{Value} \\
\midrule
\endhead
$\varepsilon$ & $e^\pi - \pi - 20$ &
$-0.00090002081052423273355702\ldots$ \\[4pt]
$c_1$ & $N_c^2/\mathcal{D} = 9/47$ &
$0.19148936170212765957446809\ldots$ \\[4pt]
$c_2$ & $(N_{\mathrm{eff}}\!-\!2N_{\mathrm{base}})/N_{\mathrm{base}}^3 = 5/64$ &
$0.07812500000000000000000000\ldots$ \\[4pt]
$c_3$ & $N_{\mathrm{base}}/(N_c\!\cdot\!\mathcal{D}) = 4/141$ &
$0.02836879432624113475177305\ldots$ \\[4pt]
$c_4$ & $(N_c\!\cdot\!\mathcal{D})/(b_3\!+\!N_{\mathrm{base}}) = 141/11$ &
$12.818181818181818181818182\ldots$ \\[4pt]
$\alpha^{-1}$ & (4-term result) &
$137.03599917700004140583386\ldots$ \\
\bottomrule
\end{longtable}

\noindent CODATA 2022: $\alpha^{-1} = 137.035\,999\,177(21)$. \quad Agreement: $< 0.001$~ppt.

%=======================================================================
\section*{Layer 4: Framework Integers}
%=======================================================================

\begin{longtable}{lp{6.5cm}r}
\toprule
\textbf{Symbol} & \textbf{Formula} & \textbf{Value} \\
\midrule
\endhead
$N_c$ & $\lfloor x_- \rfloor$ & $3$ \\
$N_{\mathrm{base}}$ & $2^{(D+1)/2}$ & $4$ \\
$b_3$ & $(11N_c - 2N_f)/3$, \quad $N_f = 2N_c = 6$ & $7$ \\
$N_{\mathrm{eff}}$ & $b_3 + 2N_c$ & $13$ \\
$\mathcal{D}$ & $N_c N_{\mathrm{base}}^2 - 1$ & $47$ \\
\bottomrule
\end{longtable}

%=======================================================================
\section*{Layer 5: Coupling Constants}
%=======================================================================

\begin{longtable}{lp{5.5cm}r}
\toprule
\textbf{Symbol} & \textbf{Formula} & \textbf{Value} \\
\midrule
\endhead
$\alpha_{\mathrm{tree}}$ & $1/x_+$ &
$0.0072973433901380833826792\ldots$ \\[4pt]
$\alpha$ & $1/\alpha^{-1}_{\mathrm{prec}}$ \quad (4-term; used below) &
$0.0072973525643314228253291\ldots$ \\[4pt]
$g_c$ & $\sqrt{\alpha}$ &
$0.0854245431028543624444677\ldots$ \\[4pt]
$\alpha_s(M_Z)$ & $b_3/(b_3 + 4N_{\mathrm{eff}}) = 7/59$ &
$0.1186440677966101694915254\ldots$ \\[4pt]
$\sin^2\!\theta_W$ & $N_c/N_{\mathrm{eff}} = 3/13$ &
$0.2307692307692307692307692\ldots$ \\[4pt]
$G_N$ & $1/(b_3 + N_c)^2 = 1/100$ &
$0.0100000000000000000000000\ldots$ \\
\bottomrule
\end{longtable}

%=======================================================================
\section*{Layer 6: Mass Scales}
%=======================================================================

\begin{longtable}{lp{5.5cm}r}
\toprule
\textbf{Symbol} & \textbf{Formula} & \textbf{Value} \\
\midrule
\endhead
$m_e/m_P$ & $\sqrt{2\pi}\cdot\frac{16}{3}\cdot\alpha^{11}$ &
$4.1774117506818401689118\ldots\times 10^{-23}$ \\[6pt]
$K_B = m_e$ & (lattice natural unit) & $1$ \quad ($= 0.5100$~MeV) \\[4pt]
$K_{\mathrm{genesis}}$ & $N_c \cdot K_B$ & $3$ \quad ($= 1.530$~MeV) \\
\bottomrule
\end{longtable}

%=======================================================================
\section*{Layer 7: Mass Ratios (Exact Integers)}
%=======================================================================

\begin{longtable}{lp{6cm}r}
\toprule
\textbf{Ratio} & \textbf{Formula} & \textbf{Value} \\
\midrule
\endhead
$m_\mu/m_e$ & $3b_3(b_3 + N_c) - N_c$ & $207$ \\[4pt]
$m_\tau/m_e$ & $(N_{\mathrm{eff}} + N_{\mathrm{base}})\cdot 207 - 2N_c b_3$ & $3477$ \\[4pt]
$m_p/m_e$ & $N_{\mathrm{eff}}/\alpha + T(b_3\!+\!N_c)$, \; $T(10) = 55$ &
$1836.4679893010005382758\ldots$ \\
\bottomrule
\end{longtable}

%=======================================================================
\section*{Layer 8: Higgs Sector}
%=======================================================================

\begin{longtable}{lp{5.5cm}r}
\toprule
\textbf{Symbol} & \textbf{Formula} & \textbf{Value} \\
\midrule
\endhead
$m_P$ & $K_B / (\sqrt{2\pi}\cdot\tfrac{16}{3}\cdot\alpha^{11})$ &
$2.3938267513054687\ldots \times 10^{22}$~lattice units\rlap{$^{\dagger}$} \\[4pt]
$v$ & $m_P\sqrt{2\pi}\;\alpha^8$ &
$246.07951397853684\ldots$~GeV \\[6pt]
$\lambda_H$ & $\displaystyle\frac{\sin^2\!\theta_W}{2 - \sin^2\!\theta_W} = \frac{3}{23}$ &
$0.1304347826086956521739130\ldots$ \\[8pt]
$m_H$ & $v\sqrt{6/23}$ &
$125.68607601610095\ldots$~GeV \\
\bottomrule
\end{longtable}

\noindent{\small $^{\dagger}$\textit{Unit conversion.} $m_P$ is computed in lattice units (l.u.), where $K_B = 1$. To obtain GeV values for $v$ and $m_H$: multiply $m_P$ by $K_B = 0.5100\times10^{-3}$~GeV, giving $m_P \approx 1.22085\times 10^{19}$~GeV; then $v = m_P\sqrt{2\pi}\,\alpha^8$ and $m_H = v\sqrt{6/23}$ follow in GeV. No additional input enters.}

\medskip

%=======================================================================
\section*{Derived Quantities}
%=======================================================================

\begin{longtable}{lp{5.5cm}r}
\toprule
\textbf{Symbol} & \textbf{Formula} & \textbf{Value} \\
\midrule
\endhead
$W_3$ & $\displaystyle\frac{\Gamma(1/4)^4}{4\pi^3} = \frac{\Gstar^2}{2\pi}$ &
$1.3932039296856768591842463\ldots$ \\[8pt]
$c$ & $1/\sqrt{3}$ &
$0.5773502691896257645091488\ldots$ \\[4pt]
$r_{\mathrm{eff}}$ & $\sqrt{x_+} = \sqrt{\alpha_{\mathrm{tree}}^{-1}}$ &
$11.706244976855536653552629\ldots$ \\[4pt]
$1/17$ & $1/(N_{\mathrm{eff}} + N_{\mathrm{base}})$ &
$0.0588235294117647058823529\ldots$ \\
\bottomrule
\end{longtable}

%=======================================================================
\section*{Confinement (from $x_-$)}
%=======================================================================

\begin{longtable}{lp{5.5cm}r}
\toprule
\textbf{Symbol} & \textbf{Formula} & \textbf{Value} \\
\midrule
\endhead
$u_p(x_-)$ & $I_1(x_-)/I_0(x_-)$ &
$0.8117414116844551696351218\ldots$ \\[4pt]
$\sigma(x_-)$ & $-\ln\,u_p$ &
$0.2085734480558214996554668\ldots$ \\
\bottomrule
\end{longtable}

%=======================================================================
\section*{Key Identities}
%=======================================================================

\begin{longtable}{lp{5.5cm}r}
\toprule
\textbf{Identity} & \textbf{Verification} & \textbf{Value} \\
\midrule
\endhead
$e^\pi$ & Gelfond's constant &
$23.140692632779269005729086\ldots$ \\[4pt]
$e^\pi - \pi - 20$ & $= \varepsilon$ &
$-0.00090002081052423273355702\ldots$ \\[4pt]
$4\ell$ & $= \pi$ &
$3.1415926535897932384626434\ldots$ \\[4pt]
$16\Gstar^2$ & $= x_+ + x_-$ &
$140.06013537449450488555328\ldots$ \\[4pt]
$16\Gstar^3$ & $= x_+ \cdot x_-$ &
$414.39243772270942760105875\ldots$ \\
\bottomrule
\end{longtable}

\bigskip

\begin{center}\itshape
All values computed with \texttt{mpmath} at 80-digit internal precision; 25 digits displayed.\\
Canonical source: \texttt{scripts/constants.py}.\\
Engine mirror: \texttt{engine/include/ftd/ontic.h} (\texttt{constexpr}).
\end{center}

\end{document}
